\documentclass[10pt,a4paper]{amsart}
\usepackage[margin=19mm]{geometry}
\usepackage{amsmath,amssymb,mathtools}
\usepackage[scr=rsfs,cal=euler]{mathalfa}
\usepackage{xfrac}
\usepackage[unicode,hidelinks,bookmarks=false]{hyperref}
\newcommand{\ie}{i.e.\ }
\newcommand{\cf}{cf.\ }
\newcommand{\resp}{respectively\ }
\newcommand{\Subsection}{\subsection}
\providecommand{\nobreakdash}{\nobreak\hbox{-}}
\setlength{\emergencystretch}{2em}
\multlinegap0pt
\usepackage[shortlabels]{enumitem}
\usepackage{amssymb}
\usepackage{tikz}
\usepackage{mathtools}
\newtheorem{lemma}{Lemma}
\newcommand{\Sph}{\mathbb{S}}
\newcommand{\V}{\mathbb{V}}
\title{Equilibria in non-Euclidean geometries}
\author{Z. L\'angi, S. Wang}
\date{Source: arXiv:2602.01159v2 (CC BY 4.0)}
\begin{document}
\maketitle

\noindent\emph{Source and licence.} Equilibria in non-Euclidean geometries. Version arXiv:2602.01159v2 (CC BY 4.0), distributed by arXiv under the Creative Commons Attribution 4.0 International licence, http://creativecommons.org/licenses/by/4.0/, as recorded in the arXiv licence metadata for this exact version and shown on the abstract page https://arxiv.org/abs/2602.01159v2. Full licence notice: \texttt{licenses/arxiv-2602.01159-CC-BY-4.0.txt}.\par\smallskip

\noindent\textit{Editorial scope.} Counted unit: the article's lemma lem:centeringnormed (source0.tex lines 698--701). The article's complete original proof of it is delivered with it (source0.tex lines 703--748, one \texttt{proof} environment of 46 lines). No mathematics is written, repaired or re-derived: the statement and the proof are the article's own lines, in the article's own notation, and no retained line is substituted or re-flowed.  Retained uncounted as the source-local context the proof needs: lines 371--371.  Nothing was omitted from any retained span, so the package carries no omission record.  No retained line contains a citation, so the wrapper carries no bibliography and no bibliography entry.  No retained line contains a figure environment, an image-inclusion command or drawing code, so no image is delivered and the package declares no asset dependency.  Editorial formatting only, in addition: an AMS \texttt{amsart} preamble that restates the article's own declarations and macros used by the retained lines, namely the environments \texttt{lemma} on separate counters, together with the standard packages those lines need.  The retained environments print in this document as Lemma 1 in order of appearance, whereas the article prints the same items with the numbering of its own full document.  The complete native source of the article as distributed in the arXiv e-print for this exact version is bundled unchanged as \texttt{sources/193765/source0.tex}.
\subsection{The case $\V^3 = \Sph^3$}\label{subsec:spherical}

\begin{lemma}\label{lem:centeringnormed}
There exist constants $0 < c_1 < c_2 < 1$, $0 < d_0 < 1$ and a function $F:(0,d_0) \to [c_1,c_2]$ such that for any $d \in (0,d_0)$,
$c(K(F(d),d))=o$.
\end{lemma}

\begin{proof}
We start the proof with the observation that as $K(c,d)$ is axially symmetric to the $x_3$-axis, its centroid is of the form $c(K(c,d))=(0,0,z(c,d))$. Thus, $c(K)=o$ if and only if the first moment of $K$ with respect to the plane $H_3$ with equation $\{ x_3 = 0 \}$ is $0$. We denote this first moment by $M_3(K)$.

Changing to spherical coordinates, up to a constant factor independent of $c,d$ and setting $x=(x_1,x_2,x_3)$, $M(K_m(c,d))$ can be written as
\begin{multline*}
M_3(K(c,d)) = \int_{x \in K} x_3  \, d \lambda_3(x) =\\
= 
\int_{0}^{2\pi}\int_{-\frac{\pi}{2}}^{\frac{\pi}{2}}\int_{0}^{R \rho_M(1 + d \rho_c)} r^3 \sin \theta \cos \theta \,dr\,d\theta\,d\varphi,
\end{multline*}
implying that
\begin{equation}\label{eq:firstmomentnormed}
M_3(K(c,d)) = \int_{0}^{2\pi}\int_{-\frac{\pi}{2}}^{\frac{\pi}{2}} \frac{1}{4} \left( R \rho_M(1 + d \rho_c) \right)^3 \sin \theta \cos \theta \,d\theta\,d\varphi .
\end{equation}

We note that as $M$ is $o$-symmetric, $M_3(K(c,0))=0$ for all $0 < c \leq 1$. Furthermore, from this it also follows that
\begin{equation}\label{eq:firstmomentnormedmod}
M_3(K(c,d)) = d \int_{0}^{2\pi}\int_{-\frac{\pi}{2}}^{\frac{\pi}{2}}  R^4 \rho_M^4 \left( \rho_c + \frac{3}{2} d \rho_c^2 + d^2 \rho_c^3 + \frac{d^3}{4} \rho_c^4\right) \sin \theta \cos \theta \,d\theta\,d\varphi .
\end{equation}

Observe that the integrand in (\ref{eq:firstmomentnormed}), as well as its partial derivative with respect to $d$, is continuous on the domain $0 < c \leq 1$ and $0 \leq d < 1$. Hence, the same holds for $M_3(K(c,d))$.
Note that $\rho_1(\theta, \varphi) = \sin(\theta)$. Thus,
\begin{multline*}
M_3(K(1,d)) =\\
= d \int_{0}^{2\pi}\int_{-\frac{\pi}{2}}^{\frac{\pi}{2}}  R^4 \rho_M^4 \left( \sin^2 \theta + \frac{3}{2} d \sin^3 \theta + d^2 \sin^4 \theta + \frac{d^3}{4} \sin^5 \theta\right) \cos \theta \,d\theta\,d\varphi
\end{multline*}
For any $-\frac{\pi}{2} \leq \theta \leq \frac{\pi}{2}$, set $\hat{\rho}_M(\theta) = \int_{0}^{2\pi} \rho_M^4(\theta,\varphi) \, d \varphi$. Then, as $M$ is $o$-symmetric, we have that $\hat{\rho}_M(\theta)$ is an even function. Thus, as the function $\theta \mapsto \sin^k \theta$ is even if $k$ is even and odd if $k$ is odd, it follows that
\begin{equation}\label{eq:firstmomentc1_2normed}
M_3(K_m(1,d)) = d \int_{-\frac{\pi}{2}}^{\frac{\pi}{2}}  2 R^4 \hat{\rho}_M \left( 2 \sin^2 \theta  + d^2 \sin^4 \theta \right) \cos \theta \,d\theta,
\end{equation}
which is positive for all values $0<d \leq 1$. In particular, $M_3(K(c,d))$ is positive in a neighborhood of the point $(c,d)=(1,0)$.

Next, we consider the value of $M_3(K(c,d))$ near the point $(c,d)=(0,0)$. Recall that $M_3(K(c,0))=0$ for any $0 < c \leq 1$. Thus, to show that $M_3(K(c,d))$ is negative for sufficiently small values of $c$ and $d$, it is sufficient to show that
\[
\bar{M}(c,d)= \int_{0}^{2\pi}\int_{-\frac{\pi}{2}}^{\frac{\pi}{2}}  \rho_M^4 \rho_c \sin \theta \cos \theta \,d\theta\,d\varphi < 0
\]
in such a region. Now, using the notation $\rho_0$ introduced in Subsection~\ref{subsec:spherical} and applying Lebesgue's dominated convergence theorem, we obtain that
\[
\lim_{c \to 0^+} \bar{M}(c,d) = \int_{0}^{2\pi}\int_{-\frac{\pi}{2}}^{\frac{\pi}{2}}  \rho_M^4 \rho_0 \sin \theta \cos \theta \,d\theta\,d\varphi .
\]
Now, since $\rho_M$ is independent of $\varphi$, it follows that
\[
\lim_{c \to 0^+} \bar{M}(c,d) = - \int_{-\frac{\pi}{2}}^{\frac{\pi}{2}} \rho_M^4 \frac{8 \theta \pi^2}{\pi^2+4 \theta^2} \sin \theta \cos \theta \,d\theta,
\]
which, as the integrand is an even function, is negative.
Combining this result with our investigation near the point $(c,d)=(1,0)$, the assertion follows.
\end{proof}

\end{document}
